\subsection{Definition of the Brauer Group}

We denote by \(\mathrm{Br}(\mathrm{K})\) the set of elements of \(M_{K}\) of the form {[}A{]}, where A is a central simple algebra of finite degree over K.

PROPOSITION 2. --- The set Br(K) is the set of invertible elements of \(\mathcal{M}_{\mathrm{K}}\) , where the inverse of an element {[}A{]} of Br(K) is {[}A°{]}. Consequently, the law of composition on the monoid \(M_{K}\) endows Br(K) with the structure of an abelian group.

Let A be a central simple algebra of finite degree over K. The algebra \(A^{o}\) is central simple and of finite degree over K (VIII, p.~251). The algebra \(A \otimes_{K} A^{o}\) is isomorphic to a matrix algebra \(\mathbf{M}_{n}(K)\) , where \(n^{2} = \dim(A)\) (VIII, p.~252, Theorem 1). We therefore have \([A][A^{o}] = [A \otimes_{K} A^{o}] = [K]\) (Lemma 1), which proves that {[}A{]} is invertible, with inverse \([A^{o}]\) .

Conversely, let A be a K-algebra of finite degree. If {[}A{]} is invertible in \(M_{K}\) , then there exists a K-algebra B of finite degree such that \([A][B]=[K]\) ; by formula (2) and Lemma 1, this means that the K-algebra \(A \otimes_{K} B\) is isomorphic to a matrix algebra \(\mathbf{M}_{n}(\mathrm{K})\) with \(n \geqslant 1\) . By Remark 1 of VIII, p.~251, the algebra A is then central simple.

DEFINITION 1. --- The abelian group Br(K) is called the Brauer group of the field K.

Lemma 2. --- Let I and J be finite sets, k be a commutative ring, and A and B be k-algebras. Denote by \(\mathbf{M}_{\mathrm{I}}(\mathrm{A})\) the k-algebra of square matrices of type (I,I) with entries in A, and define the k-algebras \(\mathbf{M}_{\mathrm{J}}(\mathrm{B})\) and \(\mathbf{M}_{\mathrm{I}\times\mathrm{J}}(\mathrm{A}\otimes_{\mathrm{K}}\mathrm{B})\) likewise. There exists a unique k-algebra isomorphism

\[
\varphi : \mathbf {M} _ {\mathrm{I}} (\mathrm{A}) \otimes_ {k} \mathbf {M} _ {\mathrm{J}} (\mathrm{B}) \longrightarrow \mathbf {M} _ {\mathrm{I} \times \mathrm{J}} (\mathrm{A} \otimes_ {k} \mathrm{B})
\]

such that \(\varphi\big((a_{ii'})\otimes(b_{jj'})\big)\) is the matrix with entry of index \(((i,j),(i',j'))\) equal to \(a_{ii'}\otimes b_{jj'}\) .

The existence of a k-linear bijection \(\varphi\) with the property stated in the lemma follows from the compatibility of the tensor product with direct sums (II, §3, No.~7, p.~255, Proposition 7). The fact that \(\varphi\) is an algebra homomorphism follows from the definition of a product matrix.

PROPOSITION 3. --- Let A and B be central simple K-algebras of finite degree. The following properties are equivalent:

\begin{enumerate}
\def\labelenumi{(\roman{enumi})}
\item
  We have {[}A{]} = {[}B{]} in the Brauer group Br(K).
\item
  There exists an integer \(t \geqslant 1\) such that the K-algebra \(\mathrm{A} \otimes_{\mathrm{K}} \mathrm{B}^{\circ}\) is isomorphic to the matrix algebra \(\mathbf{M}_t(\mathrm{K})\) .
\item
  There exist strictly positive integers \(r\) and \(s\) such that the K-algebras \(\mathrm{A} \otimes_{\mathrm{K}} \mathbf{M}_r(\mathrm{K})\) and \(\mathrm{B} \otimes_{\mathrm{K}} \mathbf{M}_s(\mathrm{K})\) are isomorphic.
\item
  There exist a field D containing K and integers \(m \geqslant 1\) and \(n \geqslant 1\) such that A is isomorphic to \(\mathbf{M}_{m}(\mathrm{D})\) and B to \(\mathbf{M}_{n}(\mathrm{D})\) .
\item
  The K-algebras A and B are Morita equivalent.
\end{enumerate}

Suppose that we have \([A]=[B]\) . Since \([B^{o}]\) is the inverse of {[}B{]} in the Brauer group, we have \([K]=[B][B^{o}]=[A][B^{o}]=[A\otimes_{K}B^{o}]\) . By Lemma 1, there exists an integer \(t\geqslant1\) such that the algebras \(A\otimes_{K}B^{o}\) and \(\mathbf{M}_{t}(K)\) are isomorphic. So (i) implies (ii).

Suppose that (ii) holds. Since \(B^{o} \otimes_{K} B\) is isomorphic to a matrix algebra \(\mathbf{M}_{s}(K)\) with \(s \geqslant 1\) (VIII, p.~252, Theorem 1), the algebra \(A \otimes_{K} B^{o} \otimes_{K} B\) is isomorphic to \(A \otimes_{K} M_{s}(K)\) , on the one hand, and to \(\mathbf{M}_{t}(K) \otimes_{K} B\) , on the other. So property (ii) implies property (iii).

Suppose that (iii) holds. By Wedderburn's theorem (VIII, p.~120, Theorem 1), there exist integers \(m \geqslant 1\) and \(n \geqslant 1\) and fields D and D' of finite degree over K with center K such that A is isomorphic to \(\mathbf{M}_m(\mathrm{D})\) and B to \(\mathbf{M}_n(\mathrm{D}')\) . Then the algebra \(\mathrm{A} \otimes_{\mathrm{K}} \mathbf{M}_r(\mathrm{K})\) is isomorphic to \(\mathbf{M}_{mr}(\mathrm{D})\) , and the algebra \(\mathrm{B} \otimes_{\mathrm{K}} \mathbf{M}_{r'}(\mathrm{K})\) is isomorphic to \(\mathbf{M}_{nr'}(\mathrm{D}')\) (Lemma 2). By Corollary 2 of VIII, p.~121, the K-algebras D and D' are isomorphic. So (iii) implies (iv).

Suppose that (iv) holds. The algebras A and D are Morita equivalent, as are the algebras B and D (VIII, p.~114, Example 2). Therefore, A and B are Morita equivalent, so (iv) implies (v).

Finally, the implication \((v)\Rightarrow(i)\) is the definition of the set \(\mathrm{Br}(K)\) .

When the equivalent properties of the proposition hold, we say that the algebras A and B are similar.

COROLLARY. --- Let A and B be central simple algebras of finite degree over K. Then A and B are isomorphic if and only if they are similar and have the same degree.

This follows from the equivalence of properties (i) and (iv) of Proposition 3 and the fact that \(\dim_{\mathrm{K}}(\mathbf{M}_n(\mathrm{D})) = \dim_{\mathrm{K}}(\mathrm{D}) \times n^2\) .

PROPOSITION 4. --- Let \(\mathcal{H}_{\mathrm{K}}\) be the set of classes of central K-algebras of finite degree that are fields. The mapping \(\mathrm{D} \mapsto [\mathrm{D}]\) from \(\mathcal{H}_{\mathrm{K}}\) to \(\mathrm{Br}(\mathrm{K})\) is bijective.

This follows from Lemma 1 of VIII, p.~278 and Wedderburn's theorem (VIII, p.~120, Theorem 1).

